  % ======================================================================
  %  SPARC - Industrial report for CONCAST
  %  Built on the personal LaTeX template (template/doc.tex).
  %  Compile: latexmk -pdf sparc_industrial_report.tex
  % ======================================================================

  \documentclass[10pt,a4paper]{scrartcl}

  % --------------------------------------------------
  % Encoding and language
  % --------------------------------------------------
  \usepackage[T1]{fontenc}
  \usepackage[utf8]{inputenc}
  \usepackage[english]{babel}

  % --------------------------------------------------
  % Page layout
  % --------------------------------------------------
  \usepackage[
    top=18mm,
    bottom=18mm,
    left=20mm,
    right=20mm,
    includefoot
  ]{geometry}

  % --------------------------------------------------
  % Mathematics and graphics
  % --------------------------------------------------
  \usepackage{graphicx}
  \usepackage{amsmath,amssymb}
  \usepackage{bm}
  \setlength{\jot}{0.25cm}

  \usepackage{tikz}
  \usetikzlibrary{positioning,arrows.meta,shapes.geometric,patterns,calc}
  \usepackage{pgfplots}
  \pgfplotsset{compat=1.18}

  % --------------------------------------------------
  % Tables
  % --------------------------------------------------
  \usepackage{booktabs}
  \usepackage{float}
  \usepackage[font=small,labelfont=bf,justification=raggedright,singlelinecheck=true]{caption}
  \usepackage{tabularx}
  \usepackage{array}
  \usepackage{ragged2e}

  % A ragged-right fixed-width column for use in tabularx:
  %   \begin{tabularx}{\textwidth}{@{}L{3cm} X@{}}
  \newcolumntype{L}[1]{>{\RaggedRight\arraybackslash}p{#1}}
  \newcolumntype{Y}{>{\RaggedRight\arraybackslash}X}
  \newcolumntype{R}{>{\raggedleft\arraybackslash}X}

  % --------------------------------------------------
  % Typography
  % --------------------------------------------------
  \usepackage[scaled=0.95]{helvet}
  \usepackage[expansion=false]{microtype}   % expansion=false avoids font errors
  \renewcommand{\familydefault}{\sfdefault} % sans-serif throughout
  \usepackage[gen]{eurosym}
  \newcommand{\Eur}[1]{\euro\,#1}

  % --------------------------------------------------
  % Lists
  % --------------------------------------------------
  \usepackage{enumitem}
  \setlist[itemize]{
    leftmargin=*,
    topsep=2pt,
    parsep=0pt,
    partopsep=0pt,
    itemsep=3pt
  }
  \setlist[enumerate]{
    leftmargin=*,
    topsep=2pt,
    parsep=0pt,
    partopsep=0pt,
    itemsep=3pt
  }
  \setlist[description]{
    leftmargin=0pt,
    topsep=2pt,
    parsep=0pt,
    itemsep=3pt
  }

  % --------------------------------------------------
  % Colour and hyperlinks
  % --------------------------------------------------
  \usepackage[dvipsnames]{xcolor}
  \usepackage[hidelinks]{hyperref}

  \definecolor{accent}{RGB}{0,129,88}   % change this to restyle everything
  \definecolor{accentB}{RGB}{37,99,235}
  \definecolor{warn}{RGB}{178,107,0}

  % \emph produces upright accent-coloured text instead of italic.
  % Use \textit{} when you genuinely want italics.
  \renewcommand{\emph}[1]{\textcolor{accent}{#1}}

  % Boxes for key figures and practical notes
  \usepackage[most]{tcolorbox}
  \newtcolorbox{keybox}[1]{colback=accent!6,colframe=accent!70,boxrule=0.6pt,arc=2pt,
    left=6pt,right=6pt,top=4pt,bottom=4pt,fonttitle=\bfseries\small,title={#1}}
  \newtcolorbox{notebox}[1]{colback=warn!6,colframe=warn!70,boxrule=0.6pt,arc=2pt,
    left=6pt,right=6pt,top=4pt,bottom=4pt,fonttitle=\bfseries\small,title={#1}}

  % --------------------------------------------------
  % Header and footer
  % --------------------------------------------------
  \usepackage[automark,headsepline=false,footsepline=false]{scrlayer-scrpage}
  \clearpairofpagestyles

  \ifoot{V.\,G. De Matteis}
  \cfoot{\pagemark}
  \ofoot{SPARC --- Industrial report for CONCAST}

  \setlength{\footskip}{20mm}  % pushes footer toward the page edge
  \pagestyle{scrheadings}

  % --------------------------------------------------
  % Spacing
  % --------------------------------------------------
  \setlength{\parindent}{0pt}
  \setlength{\parskip}{0.45em}

  % --------------------------------------------------
  % Section formatting
  % --------------------------------------------------
  \RedeclareSectionCommand[
    beforeskip=1.4ex plus 0.2ex minus 0.1ex,
    afterskip=0.6ex
  ]{section}
  \RedeclareSectionCommand[
    beforeskip=1.0ex plus 0.2ex minus 0.1ex,
    afterskip=0.4ex
  ]{subsection}
  \RedeclareSectionCommand[
    beforeskip=0.7ex plus 0.2ex minus 0.1ex,
    afterskip=0.3ex
  ]{subsubsection}

  \setkomafont{section}{\large\bfseries\color{accent}}
  \setkomafont{subsection}{\normalsize\bfseries\color{accent}}
  \setkomafont{subsubsection}{\normalsize\bfseries}

  % A section command that adds a faint rule beneath the title.
  % Start a new page if less than #1 is left (keeps headings with their text)
  \makeatletter
  \newcommand{\NeedSpace}[1]{\par\begingroup\dimen@=\pagegoal\advance\dimen@-\pagetotal
    \ifdim\dimen@<#1\relax\newpage\fi\endgroup}
  \makeatother
  \newcommand{\RuleSection}[1]{%
    \NeedSpace{10\baselineskip}%
    \section{#1}%
    \vspace{0.2em}{\color{accent!35}\hrule height 0.4pt}\vspace{0.3em}%
  }

  % Shared style of the charts
  \pgfplotsset{sparc/.style={
    width=\textwidth, height=5.2cm, grid=major, grid style={gray!20},
    tick label style={font=\scriptsize}, label style={font=\footnotesize},
    legend style={font=\scriptsize, draw=none, fill=none}, axis lines*=left}}

  % ======================================================================
  \begin{document}
  % ======================================================================

  % --------------------------------------------------
  % Title block
  % --------------------------------------------------
  \begin{center}
    {\Huge\bfseries\color{accent} SPARC: energy-aware production scheduling}\\[0.25cm]
    {\large Industrial report --- shift optimisation, cost simulation and cold-room control at CONCAST}\\[0.2cm]
    {\small Prepared for: CONCAST plant and energy management \quad\textbullet\quad Version 1.0 \quad\textbullet\quad October 2026}\\[0.1cm]
    {\small Victor Gabriele De Matteis}\\[0.2cm]
    {\color{accent!35}\hrule height 0.5pt}
  \end{center}

  \vspace{0.3em}

  % ======================================================================
  \RuleSection{Executive summary}
  \label{sec:summary}

  SPARC is a software tool, installed on a plant computer, that \emph{moves the flexible
  production batches} within the time window each batch is allowed to move, so that more of their
  electricity comes from the plant's own cogenerator and solar panels, or is bought when the grid
  price is lower. It then checks the effect on the electricity bill over any period, and it can
  simulate a cold room cooled with surplus on-site production.

  We tested SPARC on the real CONCAST data: the reference production week, hourly on-site
  generation from January 2024 to October 2025 and hourly grid prices. Each of the
  \textbf{92 weeks} was optimised with the generation and prices of that week.

  \begin{keybox}{Key figures}
    \begin{tabularx}{\linewidth}{@{}L{0.46\linewidth} R R@{}}
      & \textbf{Current plants} & \textbf{With FV3 solar plant} \\
      \midrule
      Grid electricity cost today (average per year)            & \Eur{221\,100} & \Eur{188\,500} \\
      Saving if every week is optimised with its own data         & \textbf{\Eur{4\,600}/year (2.1\%)} & \textbf{\Eur{5\,100}/year (2.7\%)} \\
      Saving if last week's optimised plan is reused             & \Eur{1\,600}/year & \Eur{2\,000}/year \\
      Grid electricity avoided (weekly optimisation)              & 18\,MWh/year & 21\,MWh/year \\
      Time to optimise one week                                    & about 2\,s & about 1\,s \\
    \end{tabularx}
  \end{keybox}

  \textbf{What this means.}
  \begin{itemize}
    \item The savings are \emph{real but modest}: only three of the six processes may move
      (\texttt{b\_past}, \texttt{b\_burr} and \texttt{b\_conf}), they use 16\% of
      the weekly electricity, and they may move by at most two hours.
    \item The optimisation is \emph{exact}: for the data it is given, SPARC finds the cheapest
      possible timetable and proves it (183 of the 184 weekly runs; the last one is within 1.5\% of
      the best).
    \item \emph{How SPARC is used matters more than the algorithm.} Reusing one optimised week on
      the following weeks keeps only a third of the saving, and in about one week out of four it
      costs more than the original plan. The saving is captured only by re-optimising every week
      with the expected generation and prices of the coming week.
    \item The new FV3 solar plant on its own lowers the grid bill by about \Eur{32\,600} per year
      with today's timetable; SPARC adds about \Eur{5\,100} on top and makes better use of the extra
      midday production.
  \end{itemize}

  \textbf{Recommendations} (details in Section~\ref{sec:next}).
  \begin{enumerate}
    \item Run a \textbf{weekly planning routine}: every week, optimise the coming week with the
      day-ahead prices and a solar/cogeneration forecast, then hand the plan to production.
    \item With production, \textbf{confirm the rules} the tool does not yet know: order between
      processes, cleaning time between batches, staff availability. Add them before going live.
    \item \textbf{Review the flexibility windows}: each extra hour of allowed movement is worth
      money (Section~\ref{sec:sensitivity}).
    \item Treat the cold-room module as a \textbf{pilot} until its room model is calibrated on
      measured temperatures.
  \end{enumerate}

  % ======================================================================
  \RuleSection{The opportunity}
  \label{sec:opportunity}

  \subsection{The plant's electricity today}

  Table~\ref{tab:balance} summarises an average week with the reference production timetable.
  Production needs about 44.5\,MWh per week. The cogenerator and the solar panels produce
  31.9\,MWh, but only 22.8\,MWh of it is used by production: the rest (9.1\,MWh) is produced when
  production does not need it, typically at night and at weekends. About half of the electricity is
  therefore bought from the grid while, at other times, on-site production goes unused.

  \begin{table}[H]
  \small
  \centering
  \renewcommand{\arraystretch}{1.15}
  \begin{tabularx}{\textwidth}{@{}Y R R@{}}
  \toprule
  \textbf{Average week (Jan 2024 -- Oct 2025)} & \textbf{Current plants} & \textbf{With FV3} \\
  \midrule
  Electricity used by production            & 44.5\,MWh & 44.5\,MWh \\
  Produced by the cogenerator               & 28.5\,MWh & 28.5\,MWh \\
  Produced by the solar panels              & 3.4\,MWh  & 10.9\,MWh \\
  Bought from the grid                      & 21.7\,MWh (49\%) & 18.2\,MWh (41\%) \\
  On-site production left unused            & 9.1\,MWh  & 13.1\,MWh \\
  Cost of the electricity bought            & \Eur{4\,250} & \Eur{3\,630} \\
  \bottomrule
  \end{tabularx}
  \caption{Weekly electricity balance with the current production timetable. Unused production is
  not credited (no sale to the grid is assumed).}
  \label{tab:balance}
  \end{table}

  \subsection{What can move}

  The production plan is a typical week. Each process has a constant power while it runs and a
  number of \emph{runs} (continuous blocks of activity) in the week. Three processes are fixed; the
  other three may start up to two hours earlier or later (Table~\ref{tab:processes}). \texttt{b\_past} alone
  carries 92\% of the flexible energy: it is the process that matters.

  \begin{table}[H]
  \small
  \centering
  \renewcommand{\arraystretch}{1.15}
  \begin{tabularx}{\textwidth}{@{}L{2.2cm} R R R R L{4.2cm}@{}}
  \toprule
  \textbf{Process} & \textbf{Power} & \textbf{Hours/week} & \textbf{Runs/week} & \textbf{Energy/week} & \textbf{Allowed movement} \\
  \midrule
  \texttt{s\_finit} & 100\,kW & 119.0 & 3 & 11.9\,MWh & fixed \\
  \texttt{s\_torre} & 130\,kW & 129.5 & 2 & 16.8\,MWh & fixed \\
  \texttt{s\_hiro}  & 134\,kW & 65.3  & 3 & 8.7\,MWh  & fixed \\
  \texttt{b\_past}  & 278\,kW & 23.3  & 5 & 6.5\,MWh  & up to 2\,h earlier or later \\
  \texttt{b\_burr}  & 5\,kW   & 25.0  & 9 & 0.1\,MWh  & up to 2\,h earlier or later \\
  \texttt{b\_conf}  & 19\,kW  & 24.8  & 8 & 0.5\,MWh  & up to 2\,h earlier or later \\
  \midrule
  \textbf{Total} & & & 30 & 44.5\,MWh & 7.1\,MWh flexible (16\%) \\
  \bottomrule
  \end{tabularx}
  \caption{Processes of the reference week (timetable of 22--28 September 2025, file
  \texttt{data/production\_timetable.csv}; powers and movement from \texttt{data/production.yaml}).}
  \label{tab:processes}
  \end{table}

  Figure~\ref{fig:idea} shows the principle on one batch: moving it two hours later puts most of
  it under the midday solar production, so less electricity is bought.

  \begin{figure}[H]
  \centering
  \begin{tikzpicture}
    \begin{axis}[sparc, height=4.4cm, width=0.82\textwidth, xmin=0, xmax=24, ymin=0, ymax=430,
        xlabel={Hour of the day}, ylabel={kW}, xtick={0,4,...,24},
        legend pos=north west, legend columns=3]
      \addplot[fill=yellow!35, draw=orange!70!black, smooth] coordinates
        {(6,0) (8,40) (10,180) (12,300) (13,310) (14,280) (16,150) (18,30) (19,0)};
      \addlegendentry{Solar production}
      \addplot[fill=gray!40, fill opacity=0.6, draw=gray, const plot] coordinates {(7.5,0) (7.5,278) (12,278) (12,0)};
      \addlegendentry{Batch, original time}
      \addplot[pattern=north east lines, pattern color=accentB, draw=accentB, thick, const plot]
        coordinates {(9.5,0) (9.5,278) (14,278) (14,0)};
      \addlegendentry{Batch, 2\,h later}
    \end{axis}
  \end{tikzpicture}
  \caption{Illustration (not real data): a 278\,kW \texttt{b\_past} batch moved two hours later
  overlaps much more with solar production.}
  \label{fig:idea}
  \end{figure}

  % ======================================================================
  \RuleSection{What SPARC does}
  \label{sec:tool}

  SPARC has three functions, used through a web browser:

  \begin{description}
    \item[] \textbf{Shift optimisation} --- for a production plan and a \emph{reference week}, it
      chooses the start time of every flexible run so that the electricity bought from the grid in
      that week costs as little as possible. The result is saved as a new production plan.
    \item[] \textbf{Cost simulation} --- for any period (a week, a month, a season), it calculates
      the electricity bought and its cost for one plan, or for two plans side by side. This is how
      an optimised plan is checked before it is adopted.
    \item[] \textbf{Cold-room simulation} (pilot) --- it simulates a cold room kept within its
      temperature range, cooling first with on-site production and pre-cooling when there is
      surplus (Section~\ref{sec:coldroom}).
  \end{description}

  \begin{figure}[H]
  \centering
  \begin{tikzpicture}[
      node distance=5mm and 5.5mm,
      box/.style={draw=accent!70, fill=accent!6, rounded corners=2pt, align=center,
                  font=\footnotesize, minimum height=11mm, text width=27mm},
      outbox/.style={box, draw=accentB!70, fill=accentB!6},
      arr/.style={-{Stealth[length=2mm]}, thick, gray!70!black}]
    \node[box] (data) {Hourly data\\generation, prices,\\temperature};
    \node[box, below=of data] (plan) {Production plan\\timetable +\\allowed movement};
    \node[box, right=12mm of $(data.east)!0.5!(plan.east)$, anchor=west] (opt) {Shift optimisation\\on a reference\\week};
    \node[outbox, right=of opt] (new) {Optimised\\production plan};
    \node[box, right=of new] (sim) {Cost simulation\\over a month\\or a season};
    \node[outbox, right=of sim, text width=20mm] (dec) {Decision:\\adopt the plan};
    \node[box, below=9mm of sim, draw=gray, fill=gray!8] (cold) {Cold-room\\simulation (pilot)};
    \draw[arr] (data.east) -- (opt.west);
    \draw[arr] (plan.east) -- (opt.west);
    \draw[arr] (opt) -- (new);
    \draw[arr] (new) -- (sim);
    \draw[arr] (sim) -- (dec);
    \draw[arr, dashed] (sim) -- node[right=1mm, font=\scriptsize, align=left] {unused\\production} (cold);
  \end{tikzpicture}
  \caption{How SPARC is used. Data are entered once; optimisation and checks take seconds.}
  \label{fig:workflow}
  \end{figure}

  \subsection{How the optimisation decides}

  The week is divided into 5-minute steps (2\,016 per week). For every flexible run SPARC
  considers every allowed start time --- for a run that may move two hours either way, 49
  possibilities --- and picks one per run so that the \emph{cost of the electricity bought from the
  grid in the week} is the lowest. At each 5-minute step, production first uses the on-site
  generation; only what is missing is bought, at that hour's price.

  The schedule always respects these rules:
  \begin{itemize}
    \item a run moves only within its allowed window (fixed runs never move);
    \item a run keeps its length and power;
    \item two runs of the same process never overlap (one machine cannot do two batches at once).
  \end{itemize}

  Technically this is a \emph{mixed-integer linear programme} (Appendix~\ref{app:model}) solved
  with the open-source solver HiGHS. For one week it has about 1\,100 yes/no decisions and is solved
  to proven optimality in 1--2 seconds; a 20-second limit bounds the wait in rare hard cases, and
  the page shows the progress of the calculation while it runs.

  \begin{notebox}{What the optimisation does not know (yet)}
    \begin{itemize}
      \item Order between processes: if one batch must finish before another can start, this is not
        enforced. Each process is moved on its own.
      \item Cleaning or changeover time between two runs of the same process: runs may end up back
        to back.
      \item Staff, shifts of people, deliveries and storage limits.
      \item Value of unused production (sale to the grid), power peaks and capacity charges, fuel
        cost of the cogenerator: only the energy bought from the grid is counted.
    \end{itemize}
    A proposed plan must therefore be reviewed by production planning before it is used.
  \end{notebox}

  % ======================================================================
  \RuleSection{Results on CONCAST data}
  \label{sec:results}

  \subsection{How the results were obtained}

  Every full week from Monday 1 January 2024 to Sunday 5 October 2025 (92 weeks) was optimised
  with the generation and prices of that week, and compared with the original timetable. This was
  done twice: with the current plants (\texttt{generation.csv}) and with the additional FV3 solar
  plant (\texttt{generation\_fv3.csv}). The same reference timetable was used for every week. The
  calculation uses exactly the code of the tool (Appendix~\ref{app:method}).

  \subsection{Savings}
  \label{sec:savings}

  \begin{table}[H]
  \small
  \centering
  \renewcommand{\arraystretch}{1.15}
  \begin{tabularx}{\textwidth}{@{}L{0.5\textwidth} R R@{}}
  \toprule
  \textbf{92 weeks, every week optimised with its own data} & \textbf{Current plants} & \textbf{With FV3} \\
  \midrule
  Grid cost, original timetable          & \Eur{391\,207} & \Eur{333\,499} \\
  Grid cost, optimised timetables        & \Eur{383\,014} & \Eur{324\,449} \\
  \textbf{Saving}                         & \textbf{\Eur{8\,193} (2.1\%)} & \textbf{\Eur{9\,050} (2.7\%)} \\
  Saving per week: lowest / typical / highest & \Eur{28} / \Eur{80} / \Eur{215} & \Eur{33} / \Eur{92} / \Eur{251} \\
  Saving per year (52 weeks)              & \Eur{4\,630} & \Eur{5\,120} \\
  Grid electricity avoided per year       & 18.4\,MWh & 21.4\,MWh \\
  CO$_2$ avoided per year (0.5\,kg/kWh)   & 9\,t & 11\,t \\
  Flexible runs moved (typical week)     & 21 of 22 & 21 of 22 \\
  \bottomrule
  \end{tabularx}
  \caption{Savings when each week is optimised knowing that week's generation and prices. This is
  the most that can be obtained with the current flexibility windows.}
  \label{tab:savings}
  \end{table}

  The saving is positive every week, but it varies (Figure~\ref{fig:weekly}): it depends on how
  much on-site production and how much price difference fall within two hours of the flexible
  runs. In percentage terms it is highest in spring and summer, when solar production is large
  (Table~\ref{tab:seasons}); in winter the bill is high but the flexible runs have little
  solar production to move into. The grid electricity avoided is exactly on-site production that
  was previously unused: SPARC does not reduce consumption, it uses the plant's own production
  better.

  \begin{figure}[H]
  \centering
  \begin{tikzpicture}
    \begin{axis}[sparc, ymin=0, xmin=-1, xmax=92, ylabel={Saving per week (\euro)},
        xtick={0,13,26,39,52,65,78,91},
        xticklabels={Jan 24,Apr 24,Jul 24,Oct 24,Jan 25,Apr 25,Jul 25,Oct 25},
        legend pos=north west, legend columns=2]
      \addplot[accent, thick, mark=none] coordinates {(0,51) (1,53) (2,70) (3,61) (4,150) (5,141) (6,83) (7,78) (8,109) (9,140) (10,38) (11,39) (12,38) (13,67) (14,166) (15,63) (16,93) (17,33) (18,73) (19,123) (20,51) (21,72) (22,80) (23,134) (24,53) (25,45) (26,148) (27,67) (28,126) (29,88) (30,81) (31,150) (32,43) (33,88) (34,66) (35,102) (36,107) (37,57) (38,89) (39,73) (40,116) (41,106) (42,93) (43,162) (44,127) (45,183) (46,35) (47,37) (48,42) (49,66) (50,53) (51,61) (52,40) (53,46) (54,79) (55,95) (56,58) (57,83) (58,40) (59,74) (60,50) (61,150) (62,114) (63,168) (64,91) (65,97) (66,158) (67,87) (68,119) (69,66) (70,137) (71,79) (72,148) (73,215) (74,114) (75,149) (76,88) (77,157) (78,65) (79,115) (80,44) (81,28) (82,95) (83,59) (84,35) (85,68) (86,100) (87,74) (88,107) (89,68) (90,97) (91,69)};
      \addlegendentry{Current plants}
      \addplot[accentB, thick, dashed, mark=none] coordinates {(0,53) (1,89) (2,70) (3,49) (4,143) (5,128) (6,115) (7,40) (8,91) (9,69) (10,71) (11,45) (12,38) (13,64) (14,158) (15,135) (16,85) (17,44) (18,61) (19,139) (20,54) (21,79) (22,99) (23,97) (24,93) (25,74) (26,180) (27,86) (28,127) (29,71) (30,57) (31,251) (32,77) (33,100) (34,116) (35,130) (36,80) (37,33) (38,145) (39,62) (40,151) (41,102) (42,114) (43,169) (44,149) (45,195) (46,35) (47,37) (48,42) (49,66) (50,53) (51,61) (52,40) (53,46) (54,79) (55,95) (56,58) (57,84) (58,40) (59,76) (60,52) (61,185) (62,117) (63,107) (64,124) (65,92) (66,94) (67,161) (68,155) (69,53) (70,203) (71,84) (72,155) (73,223) (74,124) (75,81) (76,149) (77,154) (78,121) (79,112) (80,49) (81,35) (82,100) (83,36) (84,45) (85,75) (86,111) (87,135) (88,158) (89,135) (90,111) (91,125)};
      \addlegendentry{With FV3}
    \end{axis}
  \end{tikzpicture}
  \caption{Saving of each week when it is optimised with its own generation and prices.}
  \label{fig:weekly}
  \end{figure}

  \begin{table}[H]
  \small
  \centering
  \renewcommand{\arraystretch}{1.15}
  \begin{tabularx}{\textwidth}{@{}L{3.3cm} R R R R R@{}}
  \toprule
  & & \multicolumn{2}{c}{\textbf{Current plants}} & \multicolumn{2}{c}{\textbf{With FV3}} \\
  \cmidrule(lr){3-4}\cmidrule(l){5-6}
  \textbf{Season} & \textbf{Weeks} & \textbf{Grid cost/week} & \textbf{Saving/week} & \textbf{Grid cost/week} & \textbf{Saving/week} \\
  \midrule
  Winter (Dec--Feb) & 22 & \Eur{7\,059} & \Eur{72} (1.0\%)  & \Eur{6\,525} & \Eur{71} (1.1\%) \\
  Spring (Mar--May) & 26 & \Eur{2\,731} & \Eur{101} (3.7\%) & \Eur{2\,184} & \Eur{107} (4.9\%) \\
  Summer (Jun--Aug) & 26 & \Eur{3\,412} & \Eur{88} (2.6\%)  & \Eur{2\,564} & \Eur{101} (3.9\%) \\
  Autumn (Sep--Nov) & 18 & \Eur{4\,233} & \Eur{95} (2.2\%)  & \Eur{3\,696} & \Eur{115} (3.1\%) \\
  \bottomrule
  \end{tabularx}
  \caption{Savings by season (weeks grouped by the month in which they start).}
  \label{tab:seasons}
  \end{table}

  \subsection{The way SPARC is used decides the saving}
  \label{sec:operation}

  The figures above assume that each week is optimised knowing its own generation and prices. In
  practice the plan is prepared before the week starts. We therefore tested two simple ways of
  working that need no forecast (Table~\ref{tab:strategies}).

  \begin{table}[H]
  \small
  \centering
  \renewcommand{\arraystretch}{1.15}
  \begin{tabularx}{\textwidth}{@{}L{0.42\textwidth} R R R@{}}
  \toprule
  \textbf{Way of working (current plants)} & \textbf{Saving per year} & \textbf{Share of the potential} & \textbf{Weeks costing more than today} \\
  \midrule
  Optimise each week with its own data (upper bound)       & \Eur{4\,630} & 100\% & none \\
  Reuse last week's optimised plan                          & \Eur{1\,600} & 34\%  & 22 of 91 \\
  One optimised plan per season, reused all season          & \Eur{1\,070} & 23\%  & 27 of 92 \\
  \midrule
  \textit{With FV3:} optimise each week / reuse last week's plan & \Eur{5\,120} / \Eur{2\,000} & 100\% / 38\% & none / 17 of 91 \\
  \bottomrule
  \end{tabularx}
  \caption{Saving actually obtained depending on how the optimised plan is prepared. The seasonal
  plans were optimised on a typical 2024 week of each season.}
  \label{tab:strategies}
  \end{table}

  The reason is simple: the shifts move by at most two hours, so the best position depends on the
  hour-by-hour pattern of the specific week --- whether the cogenerator runs, whether the morning is
  sunny, when prices peak. A plan tuned on another week moves batches into hours that may be worse
  in the week where it is applied. The example in the tool shows the same effect: the plan
  optimised on the week of 15~July~2024 saves 5.2\% on that week but less than 1\% over the whole of
  July.

  \begin{keybox}{Practical consequence}
    To capture the saving, SPARC should be run \textbf{every week, for the coming week}, with the
    best available expectation of that week: day-ahead electricity prices (known the day before),
    a solar production forecast and the planned cogenerator operation. With perfect information the
    potential is \Eur{4\,600}--\Eur{5\,100} per year; a good forecast should recover most of it,
    which is to be verified in the pilot (Section~\ref{sec:next}).
  \end{keybox}

  \subsection{The value of more flexibility}
  \label{sec:sensitivity}

  The two-hour window is the main limit. We repeated the optimisation on one week in four (23
  weeks, current plants), giving the three flexible processes a window of one to four hours each
  way (Table~\ref{tab:window}).

  \begin{table}[H]
  \small
  \centering
  \renewcommand{\arraystretch}{1.15}
  \begin{tabularx}{\textwidth}{@{}Y R R R@{}}
  \toprule
  \textbf{Allowed movement of \texttt{b\_past}, \texttt{b\_burr}, \texttt{b\_conf}} & \textbf{Saving} & \textbf{Saving per year} & \textbf{Longest solve} \\
  \midrule
  Up to 1 hour earlier or later      & \Eur{1\,133} (1.1\%) & \Eur{2\,560} & 1.5\,s \\
  Up to 2 hours (today's setting)   & \Eur{1\,884} (1.8\%) & \Eur{4\,260} & 7\,s \\
  Up to 3 hours                      & \Eur{2\,552} (2.5\%) & \Eur{5\,770} & 18\,s \\
  Up to 4 hours                      & \Eur{3\,108} (3.0\%)$^{*}$ & \Eur{7\,030}$^{*}$ & 20\,s (limit) \\
  \bottomrule
  \end{tabularx}
  \caption{Saving (each week optimised with its own data) against the width of the window, on 23
  weeks spread over the 92. $^{*}$Some weeks reached the 20-second limit: the true saving is
  slightly higher.}
  \label{tab:window}
  \end{table}

  Each extra hour of allowed movement adds roughly \Eur{1\,300}--\Eur{1\,700} per year: three hours
  each way would give about 35\% more saving than today, four hours about 65\% more. (On this
  sample of weeks, today's setting gives \Eur{4\,260} per year instead of the \Eur{4\,630} of the
  full 92 weeks; the comparison between rows is what matters.) Wider windows make the calculation
  longer: with four hours some weeks reach the 20-second limit of the web page, which should then be
  raised.

  \subsection{Comparison with the earlier method}

  SPARC also contains the first method developed in the project (a least-squares fit followed by a
  local search). On the CONCAST data it produced, \emph{in every one of the 92 weeks}, plans in
  which two batches of the same process overlap --- for example \texttt{b\_burr} running two batches
  at the same time on Monday. Such plans cannot be executed, and their apparent savings are not
  real. The new method never does this and is proven optimal. The earlier method is therefore hidden
  from users and kept only for developers' comparisons.

  % ======================================================================
  \RuleSection{Cold-room module (pilot)}
  \label{sec:coldroom}

  The cold-room module simulates one cold room over a period, in 5-minute steps:
  \begin{itemize}
    \item the room gains heat through its walls, in proportion to the insulation (U-value), the
      surface and the difference with the outside temperature (hourly measurements, Trento);
    \item a controller (proportional--derivative) sets the cooling power to keep the room near its
      target and inside the allowed range; a minimum runtime avoids frequent on/off cycles;
    \item cooling uses on-site production first; when there is a surplus, the room is cooled a little
      further (``stores cold''); grid electricity is used only if allowed;
    \item it can be linked to a cost simulation, so that the cold room only uses the production left
      unused by the processes.
  \end{itemize}

  It reports the room temperature, the energy used and how much of it came from the grid. It does
  not consider the price of electricity, and the conversion from heat to temperature uses a fixed
  calibration coefficient, not the measured behaviour of a real room.

  \begin{notebox}{Status}
    Useful to explore ideas (for example: how much of the 9\,MWh/week of unused production a cold
    room could absorb). Before its figures are used for decisions, the room model must be calibrated
    against logged temperatures and electricity use of a real CONCAST cold room.
  \end{notebox}

  % ======================================================================
  \RuleSection{Installing and running SPARC}
  \label{sec:operating}

  \subsection{Installation}

  SPARC runs on one computer with \textbf{Docker Desktop} (Windows, macOS or Linux; its licence is
  free only below a certain company size, so check it with IT) or, on a Linux server, with the free
  Docker Engine. The software folder contains two launchers: \texttt{start} (double-click) installs
  everything on the first run, starts SPARC and opens the browser; \texttt{stop} stops it. No
  database, password file or command line is needed.

  \begin{table}[H]
  \small
  \centering
  \renewcommand{\arraystretch}{1.15}
  \begin{tabularx}{\textwidth}{@{}Y r@{}}
  \toprule
  \textbf{Operation} & \textbf{Time measured} \\
  \midrule
  First installation (download and preparation, once)          & a few minutes, depending on the network \\
  First start with an empty database, CONCAST data loaded and examples calculated & 9\,s \\
  Normal start or restart                                         & 2\,s \\
  Optimise one week                                               & 0.5--2\,s \\
  Cost simulation of one month, two plans                         & 0.2\,s \\
  \bottomrule
  \end{tabularx}
  \caption{Times measured on a development laptop with Docker Desktop.}
  \end{table}

  On the first start SPARC already contains the CONCAST data (both generation scenarios, prices
  2024--2025, Trento temperatures, the reference week) and one calculated example of each function.
  The default login (\texttt{admin}/\texttt{sparc}) must be changed at the first use. By default
  SPARC can be opened only from the computer where it runs; a setting opens it to the plant
  network.

  \subsection{Using it}

  The browser shows two areas:
  \begin{itemize}
    \item \textbf{Results}: a home page explaining the steps, the list of optimisations, cost
      simulations and cold-room simulations with their charts, and each production week as a chart.
      Anything not yet calculated has a \emph{Calculate now} button.
    \item \textbf{Manage data}: forms to upload data and set up calculations, in the order of the
      work (1 Input data, 2 Production plans and cost simulations, 3 Shift optimisation, 4 Cold
      rooms). Each calculation has a \emph{Save and calculate} button and shows its result on the
      same page. Uploaded files are checked (missing hours, wrong columns) and errors are explained
      before anything is saved. When a setting changes, the old result is removed so that an
      outdated figure is never shown.
  \end{itemize}

  \subsection{Data needed}

  \begin{table}[H]
  \small
  \centering
  \renewcommand{\arraystretch}{1.15}
  \begin{tabularx}{\textwidth}{@{}L{3.2cm} L{2.6cm} Y@{}}
  \toprule
  \textbf{Data} & \textbf{Resolution} & \textbf{Content and source} \\
  \midrule
  On-site generation & hourly, CSV & kWh produced by the cogenerator and by the solar panels (plant metering; for planning: forecast) \\
  Electricity prices & hourly, CSV & price paid per kWh bought (supplier contract or day-ahead market price plus fixed components) \\
  Outside temperature & hourly, CSV & only for cold rooms (weather station or weather service) \\
  Production plan & 15 minutes, CSV + YAML & which process runs when in a typical week; power of each process; how far each shift may move \\
  \bottomrule
  \end{tabularx}
  \caption{Inputs. The file formats are shown on the upload pages and in the README.}
  \end{table}

  \subsection{Proposed weekly routine}

  \begin{enumerate}
    \item \textbf{Thursday/Friday}: update the production plan of next week if it differs from the
      typical week; upload the price and generation expectation for next week.
    \item Create a shift optimisation for next week and press \emph{Save and optimise}
      (seconds).
    \item Production planning reviews the moved batches (list and weekly chart) against the rules
      SPARC does not know; adjust and re-run if needed.
    \item \textbf{After the week}: upload the actual generation and prices and run a cost
      simulation of the plan that was used against the original one, to track the real saving.
  \end{enumerate}

  % ======================================================================
  \RuleSection{Limitations and risks}
  \label{sec:limits}

  \begin{table}[H]
  \small
  \centering
  \renewcommand{\arraystretch}{1.2}
  \begin{tabularx}{\textwidth}{@{}L{4.3cm} Y Y@{}}
  \toprule
  \textbf{Limitation} & \textbf{Effect} & \textbf{Mitigation} \\
  \midrule
  Dependencies between processes, cleaning times, staff are not modelled &
    A proposed plan may be impossible to run as it is &
    Review by production planning; add these rules to the model (simple extension) \\
  Savings measured with perfect knowledge of each week &
    Real savings will be lower, depending on forecast quality &
    Weekly routine with forecasts; track the real saving after each week \\
  Same typical week used for all 92 weeks &
    Holidays, maintenance and changes in volume are not reflected &
    Use the actual plan of each week when available \\
  Only electricity bought is counted &
    Unused production, power peaks, capacity charges and cogenerator fuel are not valued &
    Add a value for exported energy and a peak-power charge if relevant to the contract \\
  Power of each process is constant while it runs &
    Real power profiles within a batch are smoothed &
    Use measured profiles per process if available \\
  Cold-room model not calibrated &
    Temperatures and energies are indicative &
    Calibrate on logged data before use \\
  One computer, local database &
    Data are lost if the computer is lost &
    Regular backup of the SPARC data volume \\
  \bottomrule
  \end{tabularx}
  \caption{Main limitations of the current version.}
  \end{table}

  % ======================================================================
  \newpage
  \RuleSection{Recommendations and next steps}
  \label{sec:next}

  The actions are in order of priority. The first two are needed before plans proposed by SPARC are
  used in production; the others increase the saving.

  \begin{table}[H]
  \small
  \centering
  \renewcommand{\arraystretch}{1.2}
  \begin{tabularx}{\textwidth}{@{}L{0.5cm} Y L{3.2cm} L{2.6cm}@{}}
  \toprule
  & \textbf{Action} & \textbf{Expected benefit} & \textbf{Effort} \\
  \midrule
  1 & \textbf{Pilot of the weekly routine} for 8--12 weeks in spring/summer, with day-ahead prices
      and a solar forecast; record the plan used and the real saving each week. &
      Measures the share of the \Eur{4\,600}--\Eur{5\,100}/year that is reachable &
      Low: planner time, data feeds \\
  2 & \textbf{Workshop with production} to list the rules between processes (order, cleaning, staff)
      and add them to the optimisation. &
      Plans that can be used without manual correction & Low--medium \\
  3 & \textbf{Review the movement windows} of the flexible processes, starting with \texttt{b\_past};
      consider making more processes movable. &
      See Table~\ref{tab:window} & Organisational \\
  4 & \textbf{Value the unused production} (about 9\,MWh/week): export price, thermal storage or
      cold rooms; add it to the cost if it can be sold. &
      Larger than the shift saving if a use is found & Medium \\
  5 & \textbf{Calibrate the cold-room model} on one instrumented room, then add price awareness. &
      Reliable cold-room figures & Medium \\
  6 & \textbf{Automatic data feeds} for prices and generation, and a scheduled weekly optimisation. &
      Routine without manual uploads & Medium \\
  \bottomrule
  \end{tabularx}
  \caption{Recommended actions, in order of priority.}
  \end{table}

  % ======================================================================
  \appendix
  \RuleSection{Appendix: optimisation model}
  \label{app:model}

  For engineers who want to check the method. The week is divided into steps $t=1,\dots,N$ of
  5~minutes ($N=2016$). Each run $r$ has a demand profile $\bm{d}_r$ (kWh per step) and an activity
  profile $\bm{a}_r$ (1 while it runs); $T_k$ shifts a profile by $k$ steps. Run $r$ may move by
  $k \in \{-B_r,\dots,F_r\}$ (backward and forward limits of its shift). Generation $g_t$ and price
  $\pi_t \ge 0$ are known. The binary $x_{r,k}$ is 1 when run $r$ is moved by $k$; $y_t \ge 0$ is the
  energy bought at step $t$.
  \begin{equation*}
    \begin{aligned}
      \min_{x,\,y}\quad & \sum_{t=1}^{N} \pi_t\, y_t \\
      \text{s.t.}\quad & \sum_{k=-B_r}^{F_r} x_{r,k} = 1 && \text{every run } r \\
      & \sum_{r}\sum_{k} x_{r,k}\,(T_k \bm{d}_r)_t - y_t \le g_t && \text{every step } t \\
      & \sum_{r \in p}\sum_{k} x_{r,k}\,(T_k \bm{a}_r)_t \le 1 && \text{every process } p,\ \text{step } t \\
      & x_{r,k} \in \{0,1\},\quad y_t \ge 0.
    \end{aligned}
  \end{equation*}
  The second constraint with $y_t \ge 0$ makes $y_t$ equal to the energy bought,
  $\max(0, D_t - g_t)$, at the optimum (prices are not negative); the third prevents two runs of the
  same process from overlapping. The model is solved by HiGHS through \texttt{scipy.optimize.milp}
  with a relative optimality gap of $10^{-4}$ and a 20-second limit; if no better schedule is found,
  the original one is kept. Hourly generation is interpolated linearly to 5-minute steps; prices
  are constant within the hour.

  \RuleSection{Appendix: how the figures were produced}
  \label{app:method}

  \begin{itemize}
    \item Data, in the \texttt{data} folder of the software:
      \begin{itemize}
        \item \texttt{generation.csv} and \texttt{generation\_fv3.csv}: 1 January 2024 -- 10 October 2025;
        \item \texttt{prices.csv}: 2024--2025, from 0.076 to 0.365\,\euro/kWh, average 0.175\,\euro/kWh;
        \item \texttt{production\_timetable.csv} and \texttt{production.yaml}: reference week and processes.
      \end{itemize}
    \item For each of the 92 full weeks and each generation file, the optimisation problem was built
      exactly as the tool builds it (same 5-minute grid, run detection, interpolation) and solved with
      the optimiser of the tool. Costs are those of the tool's cost function.
    \item ``Reuse last week's plan'': the moves found for week $w-1$ applied to week $w$.
      ``One plan per season'': the moves found for the 2024 week of median cost of each season,
      applied to all weeks of that season in both years.
    \item Per-year figures are the 92-week (or 91-week) totals divided by the number of weeks and
      multiplied by 52. CO$_2$ uses the tool's factor of 0.5\,kg per kWh bought.
    \item Local time is used throughout; the effect of the two daylight-saving changes per year is
      negligible for these totals.
    \item The scripts are in \texttt{report/analysis} (\texttt{benchmark.py}, \texttt{strategies.py},
      \texttt{sensitivity.py}); the technical note on the optimisation explains the method in detail.
  \end{itemize}

  \end{document}
